% Блок-схема функции eval::evaluate(rpn) — вычисление RPN на стеке
% ГОСТ 19.701-90, src/eval/evaluator.cpp
%
% Структура: главная вертикальная ось — цикл по токенам RPN.
% Подпрограммы apply_operator и apply_function раскрыты на отдельных схемах
% (соединители P и F).

\documentclass[a4paper]{article}
\usepackage[T2A]{fontenc}
\usepackage[utf8]{inputenc}
\usepackage[russian]{babel}
\usepackage[a4paper,margin=1.2cm]{geometry}
\usepackage{tikz}
\usetikzlibrary{shapes.geometric, arrows.meta, positioning, calc}

\tikzset{
  terminal/.style={
    rectangle, rounded corners=8pt,
    draw=black, line width=1pt,
    minimum width=4.4cm, minimum height=0.85cm,
    text centered, font=\footnotesize
  },
  process/.style={
    rectangle,
    draw=black, line width=1pt,
    minimum width=4.4cm, minimum height=0.85cm,
    text centered, font=\footnotesize
  },
  decision/.style={
    diamond, aspect=2.6,
    draw=black, line width=1pt,
    inner sep=1pt,
    minimum width=3.2cm, minimum height=0.85cm,
    text centered, font=\footnotesize
  },
  predefined/.style={
    rectangle, double, double distance=2pt,
    draw=black, line width=1pt,
    minimum width=4.4cm, minimum height=0.85cm,
    text centered, font=\footnotesize
  },
  connector/.style={
    circle,
    draw=black, line width=1pt,
    minimum size=0.7cm,
    text centered, font=\footnotesize\bfseries
  },
  arrow/.style={->, >=Stealth, line width=1pt}
}

\begin{document}
\begin{center}
\textbf{Блок-схема функции \texttt{evaluate(rpn)}}
\end{center}

% =========================================================
% ЧАСТЬ 1. Главный цикл по токенам RPN
% =========================================================
\begin{center}
\begin{tikzpicture}[node distance=0.6cm]

\node (start) [terminal]                  {Начало evaluate};
\node (init)  [process,  below=of start]  {st = $\varnothing$ \ (стек чисел)};
\node (loop)  [decision, below=of init]   {Есть ещё токен $t$ в rpn?};

\node (kind)  [decision, below=of loop, yshift=-0.4cm]
                                          {Тип $t$?};

% Соединители веток switch
\node (N) [connector, below left=1.3cm and 3.0cm of kind]  {N};
\node (P) [connector, below=1.3cm of kind]                 {P};
\node (F) [connector, below right=1.3cm and 3.0cm of kind] {F};

% Ветка default — ошибка
\node (errBad)[terminal, right=4.0cm of kind] {Ошибка: неожиданный токен};

% Возврат всех веток к loop через соединитель R
\node (R) [connector, right=5.5cm of loop] {R};

% Финальная проверка после цикла
\node (sz)    [decision, below=2.6cm of P]  {$|\text{st}| = 1$?};
\node (errSz) [terminal, right=2.6cm of sz] {Ошибка: некорректное выражение};
\node (ret)   [terminal, below=of sz]       {Возврат st.top()};

% --- Стрелки ---
\draw [arrow] (start) -- (init);
\draw [arrow] (init)  -- (loop);

\draw [arrow] (loop) -- node[left] {Да} (kind);
\draw [arrow] (loop.east) -- node[above, font=\footnotesize] {Нет} (R.west);
\draw [arrow] (R.south) |- (sz.east);

\draw [arrow] (kind.west)  -| node[above, pos=0.2, font=\footnotesize] {Number}   (N);
\draw [arrow] (kind.south) -- node[right, font=\footnotesize] {Operator}          (P);
\draw [arrow] (kind.east)  -| node[above, pos=0.2, font=\footnotesize] {Function} (F);

\draw [arrow] (kind.east) -- ++(2.0,0) node[midway, above, font=\footnotesize] {default} -- (errBad.west);

\draw [arrow] (sz)        -- node[left]  {Да}  (ret);
\draw [arrow] (sz.east)   -- node[above, font=\footnotesize] {Нет} (errSz);

\end{tikzpicture}
\end{center}

\bigskip

\begin{center}
\begin{tikzpicture}[node distance=0.55cm]

% --- Number ---
\node (Nin)  [connector]                 {N};
\node (Npush)[process,  below=of Nin]    {st.push($t$.value)};
\node (Nout) [connector, below=of Npush] {R};
\draw [arrow] (Nin)   -- (Npush);
\draw [arrow] (Npush) -- (Nout);

% --- Operator (apply_operator) ---
\node (Pin)  [connector, right=3.5cm of Nin] {P};
\node (Pun)  [decision,  below=of Pin]       {$t$.text = "u-"?};
\node (Pue)  [decision,  below=of Pun]       {st пуст?};
\node (PerrU)[terminal,  right=2.6cm of Pue] {Ошибка: нет операнда для u-};
\node (Pua)  [process,   below=of Pue]       {a = st.top(); st.pop()};
\node (Pur)  [process,   below=of Pua]       {r = apply\_unary("u-", a)};
\node (Pbsz) [decision,  right=4.4cm of Pun] {$|\text{st}| < 2$?};
\node (PerrB)[terminal,  right=2.6cm of Pbsz]{Ошибка: недостаточно операндов};
\node (Pbb)  [process,   below=of Pbsz]      {b = st.top(); st.pop()};
\node (Pba)  [process,   below=of Pbb]       {a = st.top(); st.pop()};
\node (Pbr)  [process,   below=of Pba]       {r = apply\_binary($t$.text, a, b)};
\node (Pchk) [predefined, below=2.0cm of Pur]{check\_finite(r)};
\node (Ppush)[process,   below=of Pchk]      {st.push(r)};
\node (Pout) [connector, below=of Ppush]     {R};

\draw [arrow] (Pin)  -- (Pun);
\draw [arrow] (Pun)  -- node[left] {Да} (Pue);
\draw [arrow] (Pun.east) -- node[above, font=\footnotesize] {Нет} (Pbsz);
\draw [arrow] (Pue.east) -- node[above, font=\footnotesize] {Да}  (PerrU);
\draw [arrow] (Pue)  -- node[left] {Нет} (Pua);
\draw [arrow] (Pua)  -- (Pur);
\draw [arrow] (Pbsz.east) -- node[above, font=\footnotesize] {Да} (PerrB);
\draw [arrow] (Pbsz) -- node[left] {Нет} (Pbb);
\draw [arrow] (Pbb)  -- (Pba);
\draw [arrow] (Pba)  -- (Pbr);
\draw [arrow] (Pur)  -- (Pchk);
\draw [arrow] (Pbr.south) |- (Pchk.east);
\draw [arrow] (Pchk) -- (Ppush);
\draw [arrow] (Ppush)-- (Pout);

% --- Function (apply_function) ---
\node (Fin)  [connector, right=4.5cm of Pbr.east, anchor=west] {F};
\node (Fe)   [decision,  below=of Fin]         {st пуст?};
\node (FerrE)[terminal,  right=2.6cm of Fe]    {Ошибка: нет аргумента};
\node (Fa)   [process,   below=of Fe]          {a = st.top(); st.pop()};
\node (Fr)   [process,   below=of Fa]          {r = functions::apply($t$.text, a)};
\node (Fchk) [predefined, below=of Fr]         {check\_finite(r)};
\node (Fpush)[process,   below=of Fchk]        {st.push(r)};
\node (Fout) [connector, below=of Fpush]       {R};

\draw [arrow] (Fin)  -- (Fe);
\draw [arrow] (Fe.east) -- node[above, font=\footnotesize] {Да} (FerrE);
\draw [arrow] (Fe)   -- node[left] {Нет} (Fa);
\draw [arrow] (Fa)   -- (Fr);
\draw [arrow] (Fr)   -- (Fchk);
\draw [arrow] (Fchk) -- (Fpush);
\draw [arrow] (Fpush)-- (Fout);

\end{tikzpicture}
\end{center}

\end{document}
